\PassOptionsToPackage{table,xcdraw}{xcolor}
\documentclass[11pt, a4paper]{article}

\usepackage[T1]{fontenc}
\usepackage[utf8]{inputenc}
\usepackage{tgpagella}
\usepackage[spanish, es-noshorthands]{babel}
\usepackage{amsmath, amssymb}
\usepackage[most]{tcolorbox}
\usepackage{tabularx}
\usepackage[margin=2cm]{geometry}
\usepackage{fancyhdr}

\pagestyle{fancy}
\fancyhf{}
\setlength{\headheight}{16pt}
\lhead{\textbf{UCB Sede Tarija} | Ing. Mecatrónica}
\rhead{IMT-221 | Protocolo de Autoevaluación}
\renewcommand{\headrulewidth}{0.4pt}
\definecolor{ucbblue}{RGB}{0, 51, 102}

\begin{document}

\begin{tcolorbox}[colback=ucbblue!5, colframe=ucbblue, title=\textbf{Protocolo de Autoevaluación Individual EC2}]
    Revise su examen comparándolo con el solucionario. Asigne un puntaje provisional y tipifique sus errores. \textbf{Nota:} Este puntaje es provisional; la corrección oficial del docente es la definitiva[cite: 9].
\end{tcolorbox}

\textbf{Nombre:} \hrulefill \quad \textbf{Códigos de Error:} \textbf{A}=Aritmético, \textbf{M}=Modelo/Enfoque, \textbf{C}=Conceptual, \textbf{S}=Signo/Convención, \textbf{U}=Unidades, \textbf{J}=Justificación, \textbf{O}=Omitido[cite: 9].

\vspace{0.3cm}
\renewcommand{\arraystretch}{1.5}
\begin{tabularx}{\textwidth}{|c|p{1cm}|p{1.5cm}|c|>{\raggedright\arraybackslash}X|>{\raggedright\arraybackslash}X|}
\hline
\rowcolor{ucbblue!20}
\textbf{Ítem} & \textbf{Max} & \textbf{Mi Nota} & \textbf{Error} & \textbf{¿Qué falló? (Causa)} & \textbf{¿Qué haré diferente?} \\ \hline
P1.1 & 3 & & & & \\ \hline
P1.2 & 3 & & & & \\ \hline
P1.3 & 4 & & & & \\ \hline
P1.4 & 3 & & & & \\ \hline
P1.5 & 2 & & & & \\ \hline
\rowcolor{gray!10} \multicolumn{6}{|l|}{\textbf{P1 Subtotal:} \rule{1cm}{0.15mm} / 15} \\ \hline
P2.1 & 4 & & & & \\ \hline
P2.2 & 4 & & & & \\ \hline
P2.3 & 4 & & & & \\ \hline
P2.4 & 4 & & & & \\ \hline
P2.5 & 4 & & & & \\ \hline
\rowcolor{gray!10} \multicolumn{6}{|l|}{\textbf{P2 Subtotal:} \rule{1cm}{0.15mm} / 20} \\ \hline
P3.1 & 4 & & & & \\ \hline
P3.2 & 4 & & & & \\ \hline
P3.3 & 2 & & & & \\ \hline
P3.4 & 2 & & & & \\ \hline
P3.5 & 3 & & & & \\ \hline
\rowcolor{gray!10} \multicolumn{6}{|l|}{\textbf{P3 Subtotal:} \rule{1cm}{0.15mm} / 15} \\ \hline
P4.1 & 3 & & & & \\ \hline
P4.2 & 5 & & & & \\ \hline
P4.3 & 4 & & & & \\ \hline
P4.4 & 4 & & & & \\ \hline
P4.5 & 4 & & & & \\ \hline
\rowcolor{gray!10} \multicolumn{6}{|l|}{\textbf{P4 Subtotal:} \rule{1cm}{0.15mm} / 20} \\ \hline
P5.1 & 4 & & & & \\ \hline
P5.2 & 4 & & & & \\ \hline
P5.3 & 3 & & & & \\ \hline
P5.4 & 5 & & & & \\ \hline
P5.5 & 5 & & & & \\ \hline
P5.6 & 3 & & & & \\ \hline
P5.7 & 6 & & & & \\ \hline
\rowcolor{gray!10} \multicolumn{6}{|l|}{\textbf{P5 Subtotal:} \rule{1cm}{0.15mm} / 30} \\ \hline
\end{tabularx}

\vspace{0.5cm}
\begin{tcolorbox}[colback=white, colframe=black, title=\textbf{Resumen y Reflexión}]
    \noindent \textbf{Puntaje Bruto Autoevaluado:} \rule{2cm}{0.15mm} / 100 \\[0.2cm]
    \noindent Dominio más fuerte: \hrulefill \\[0.2cm]
    \noindent Dominio más débil (A reforzar): \hrulefill \\[0.2cm]
    \noindent Error conceptual más grave cometido: \hrulefill \\[0.2cm]
    \noindent Pregunta o duda persistente para el docente: \hrulefill
\end{tcolorbox}

\end{document}
